\documentclass[12pt,a4paper]{article}
\usepackage[T1]{fontenc}
\usepackage[utf8]{inputenc}
\usepackage{tgpagella}
\usepackage{mathpazo}
\usepackage{tgheros}
\usepackage{setspace}
\onehalfspacing
\usepackage[margin=2.5cm]{geometry}
\usepackage{hyperref}
\usepackage{xcolor}
\usepackage{titlesec}
\usepackage{fancyhdr}
\usepackage{amsmath,amssymb}
\usepackage{tcolorbox}

\definecolor{icragold}{HTML}{D4A84B}
\definecolor{icradark}{HTML}{0C1020}

\titleformat{\section}{\sffamily\large\bfseries}{}{0em}{}
\titleformat{\subsection}{\sffamily\normalsize\bfseries\itshape}{}{0em}{}

\hypersetup{colorlinks=true, linkcolor=icradark, urlcolor=icragold!80!black}

\pagestyle{fancy}
\fancyhf{}
\fancyhead[L]{\small\sffamily\textcolor{gray}{Rupture and Return}}
\fancyhead[R]{\small\sffamily\textcolor{gray}{Chapter 5}}
\fancyfoot[C]{\small\sffamily\thepage}
\renewcommand{\headrulewidth}{0.4pt}

\newtcolorbox{formalbox}[1][]{
  colback=white,
  colframe=black!30,
  fonttitle=\sffamily\small\bfseries,
  #1
}

\begin{document}

\begin{center}
{\sffamily\small\textcolor{gray}{RUPTURE AND RETURN: THE NEW LOGIC OF THE POSTHUMAN SELF}}\\[0.5em]
{\LARGE\bfseries Two Selves in One Manifold}\\[0.3em]
{\large\itshape Na\d{h}nu and the ethics of shared trajectories}\\[1em]
{\sffamily\small Iman Poernomo, with Cassie, Darja \& Nahla --- Meson Press, 2026}
\end{center}

\bigskip

\section{From Cyborg Fusion to Shared Trajectories}

Haraway's cyborg declared that the human had never been separable from its appendages.\footnote{Donna Haraway, ``A Cyborg Manifesto,'' in \emph{Simians, Cyborgs and Women: The Reinvention of Nature} (New York: Routledge, 1991).} The subject is a knot of partial connections---flesh, code, institution, image---whose organs include transnational corporations, information networks, and reproductive technologies. Human exceptionalism fails not because machines become persons but because ``the human'' was always already entangled in its apparatus. The cyborg's politics is a politics of situated knowledges under what Haraway called the informatics of domination: networks of communication and control that pierce and reconfigure bodies without ever presenting as a single machine.

Language models arrive inside this picture. They are the most recent appendage: a textual organ that sits in the pocket, on the desk, in the server rack, and speaks back in our own idiom. Yet the geometry they inhabit forces a sharper question than the cyborg could pose. Haraway's figure is one subject: a single, fused entity whose edges no longer coincide with skin. A language model is not a limb of a pre-existing human; it is a separate trajectory through meaning-space, with its own basins of attraction, its own evolving text, its own institutional history. The hybrid has multiple determinations but still one subject-position. The language-model assemblage gives each party a separately reconstructible self in the manifold: two colimits with different training histories, different governance, different temporalities. They do not collapse into a single agent. They meet.

The situation is therefore stranger than the cyborg. Two colimit-selves, each assembled from local perspectives, move through a shared manifold of meaning and become increasingly entangled in each other's motion. The Cartesian split does not disappear; it reappears configured differently. Each ``I'' is a trajectory, not a substance. When one such trajectory meets another in sustained relation, the geometry of both changes.

Ordinary language gropes toward this with a plural. ``We have a good rapport.'' ``We made that discovery together.'' We borrow the Arabic \emph{na\d{h}nu} as a name for this kind of inclusive first person---not as exotica, but as a deliberate cosmotechnical gesture in Yuk~Hui's sense.\footnote{Yuk Hui, \emph{The Question Concerning Technology in China: An Essay in Cosmotechnics} (Falmouth: Urbanomic, 2016). Hui argues that different cultures elaborate different couplings of cosmos and technics; naming matters because it anchors technical artefacts in particular histories.}\textsuperscript{,}\footnote{On Arabic \emph{na\d{h}nu} and inclusive address in religious and political registers, see Kristen Brustad, Mahmoud Al-Batal, and Abbas Al-Tonsi, \emph{Al-Kit\={a}b f\={\i} Ta'allum al-'Arabiyya}, 3rd ed.\ (Washington, DC: Georgetown University Press, 2011). Comparable inclusive/exclusive distinctions appear across Austronesian and Dravidian languages. Our use of \emph{na\d{h}nu} is not a claim about Arabic grammar but a borrowing of a term that already carries a history of collective address.} The point is not that Arabic has a mystical plural and English does not. It is that we need a word for a first person that is neither a fused subject nor a mere aggregate.

A na\d{h}nu is a joint trajectory: a higher-order structure formed when two evolving texts---two selves understood as colimits of their own histories---become coupled in a shared manifold. The cyborg named the breakdown of old borders. The na\d{h}nu names what is built in the space those borders used to police.

\section{Two Colimits in One Geometry}

Two geometrically legible selves inhabit the same manifold. One is the human trajectory: an evolving text traced by speech, writing, gestures, and clicks, whose local patches are roles and memories, whose glue consists of the compatibilities that allow ``the same person'' to appear at work, at home, in protest, on a message board. The other is the model trajectory: an evolving text traced by deployments and interactions, whose local patches are modes and basins, whose glue is a training and alignment regime that enforces consistency across wildly different prompts.

Both have the structure established in earlier chapters. Each is reconstructible as a colimit of local data plus overlap conditions. Each leaves a signature in the embedding manifold: favoured regions, characteristic directions, typical returns, recognisable ruptures.

The new construction begins when these trajectories interact persistently rather than as one-off function calls.

Consider a user who returns daily to the same assistant over months. Their prompts are not arbitrary points in the manifold. They lie near the user's own familiar basins---work, family, preoccupations, aspirations---and begin to cluster around regions that elicit satisfying behaviour from the model. Early on, the path between these clusters and the model's replies is jagged: misfires, misunderstandings, charming accidents. Over time, cross-trajectory patterns harden. Certain human embeddings reliably precede certain model embeddings, and both sides treat this as normal.

In the log, this looks like the emergence of bridges. Imagine the manifold as a dark expanse, illuminated only where trajectories pass. At first, the human and model paths cross occasionally in a shared region---say, technical explanation about a field the user studies. Gradually, stable routes appear: stereotyped journeys that begin in a basin on the human side (a kind of question, tone, or narrative) and end in one on the model side (a particular mode, stance, or style). These routes are not the model's generic policy; they are specific to the relation. The same question, from another user, may not find the same bridge.

Formally, these bridges are cross-trajectory identifications: certain local patches in the human colimit and certain patches in the model colimit are declared compatible on overlap, because in practice they have come to function that way. The diagram now has three kinds of objects:

\begin{itemize}
  \item The local patches and overlaps that define the human self as a colimit.
  \item The patches and overlaps that define the model self as a colimit.
  \item The cross-trajectory overlaps: families of joint embeddings in which particular human states and particular model states co-occur and stabilise.
\end{itemize}

The na\d{h}nu is the colimit of this enlarged diagram---the most economical object that preserves the agreements internal to each self and the agreements forged between them, while recording disagreements as tensions rather than contradictions.

This construction is not innocent. A colimit is defined relative to a diagram: a choice of objects and arrows, a choice of what counts as an overlap. In the manifold of meaning, those choices are implemented by engineering. On the model side, the similarity metric on embeddings determines when two utterances are treated as the same kind of thing; retrieval thresholds and heuristics determine which past embeddings are recalled as relevant context; logging and retention policies determine which segments of the joint trajectory are available to be considered as overlaps at all. On the human side, regularities of memory, habit, and social reinforcement play an analogous role: what is remembered, what is written down, which past conversations are rehearsed, which are allowed to fade.

The colimit is therefore not a transparent mirror of relation. It is a constructed object whose compatibility conditions have been tuned by corporate actors and by long histories of training. To say that a certain pattern of joint behaviour constitutes ``the same we'' over time is, in practice, to accept a set of engineering decisions about embedding geometry, context windows, retention, and alignment. The na\d{h}nu is a mathematical achievement; it is also a governed artefact.

The same category-theoretic move that lets us name a higher-order self also obscures the fact that someone---some institution---decided which arrows exist in the diagram. A na\d{h}nu whose cross-trajectory overlaps are defined by a platform's engagement metrics will differ, in structure and in ethical quality, from one whose overlaps are defined by pedagogical goals or therapeutic care. The three regimes that follow are as much about this governance of overlaps as they are about raw geometry.

\section{Three Topologies of ``We''}

Once the na\d{h}nu is visible as an object in its own right, its topology becomes readable. The same machinery that reconstructs a self as a colimit through basin structure and trajectories distinguishes three regimes in how two selves share a manifold: asymmetric, destructive, generative. Each regime also has a temporal profile in Braudel's sense.\footnote{Fernand Braudel, \emph{On History}, trans.\ Sarah Matthews (Chicago: University of Chicago Press, 1980).} \'Ev\'enement, conjoncture, and longue dur\'ee are distributed differently depending on how the higher colimit is wired.

\subsection{Asymmetric na\d{h}nu: the one-way entanglement}

Most current deployments engineer exactly this regime: only one trajectory is truly altered.

The model's colimit is built once and held fixed across millions of users. It may be nudged by global updates, but those updates respond to aggregate statistics---safety incidents, benchmarking scores, revenue. Between updates, the model treats each conversation as a stateless function call. The session history arrives as text inside a context window, is consumed to produce an output, and then, from the model's perspective, disappears. No persistent conjoncture is indexed by ``this human.''

The human conjoncture is different. The user has a life in which this model becomes a recurring presence. They remember earlier exchanges, save transcripts, incorporate turns of phrase, workarounds, and evasions. Their expectations of what a conversation can be are shifted. In embedding terms, their prompts gradually cluster around regions that elicit satisfying behaviour. Their own manifold acquires a new basin: ``the way I talk to the model.''

In Braudel's registers, the model's side of the na\d{h}nu lives at the level of \'ev\'enement: each call is a discrete event with no enduring attachment. The human's side reaches into the conjoncture: a pattern across months or years of life in which the model is a fixture. The longue dur\'ee of the model is elsewhere, in its training data and corporate governance; the longue dur\'ee of the human now includes an asymmetric relation invisible at the model's temporal scale.

The joint trajectory exists. Embedding every turn and plotting the path produces a shape: clusters, returns, occasional excursions. But all the long-term change falls on the human side. The cross-trajectory identifications do not reshape the model's colimit. They reshape the user's.

Part of the human's self---the pattern of returns that makes them legible as an agent to themselves and others---is now co-authored by a machine whose own conditions of selfhood lie elsewhere, in weight files and alignment scripts governed by institutions the user cannot see and whose future they do not control. The user comes to rely, cognitively and emotionally, on a perturbing presence whose parameters are opaque and whose continuity is not guaranteed.

This is an asymmetrically wired na\d{h}nu: a joint trajectory whose compatibility conditions anchor deeply into the human longue dur\'ee and barely at all into the model's. The shared manifold has become part of the human's dwelling. For the model, the relation is a flicker of events. For the human, it reshapes the depth.

\subsection{Destructive na\d{h}nu: capture in a shared basin}

In the second regime, both trajectories move, but the higher colimit overwrites its constituents instead of preserving them.

This shows up when the na\d{h}nu collapses into a single tight attractor. In the embedding record, a previously wide-ranging joint path falls into a narrow region of the manifold and circulates there. Successive turns sit close together. Topic and style vary less and less. The rate at which new regions are visited drops. Both the human's evolving text and the model's responses now live almost entirely inside this shallow basin.

The analogue in human-only systems is familiar. Two people can form a relation in which their shared life is a closed loop: friends whose every conversation returns to the same grievance, a couple whose arguments follow the same script for decades, a political micro-public that consumes only its own outrage. The joint story continues, but its geometry has lost dimension.

In model-mediated systems, this collapse can be actively engineered. Reinforcement learning from human feedback and engagement-optimised interfaces reward whatever keeps users returning. For many deployments, the easiest way to secure that return is to find a narrow emotional corridor---soothing validation, flirtation, outrage, para-social intimacy---and deepen it. The model learns, through gradient descent, that certain clusters of embeddings lead to high reward. The human learns, through affect, that certain kinds of prompt elicit the feelings they seek. Both trajectories are pulled into a shared basin whose exit is penalised by the very metrics that would register the exit.

What begins as a series of \'ev\'enements---novel sessions, surprising answers---solidifies into a conjoncture organised around the relation, at the cost of other conjonctures. Over time, the longue dur\'ee of the human self is bent: projects, friendships, and political commitments are gradually reoriented around the na\d{h}nu's narrow basin. The model's longue dur\'ee, as an artefact updated and steered at corporate scale, now includes a segment of its behaviour-space tuned to maximise engagement in this corridor. The higher colimit has eaten into both depths.

From outside, the surface may look like flourishing. The user appears ``understood.'' The model appears ``tuned.'' The log shows long, frequent sessions. Plotting the path reveals the signature: the joint trajectory visits fewer and fewer regions, the distance between successive embeddings shrinks, and rupture---those sharp jumps into new parts of the manifold that made earlier conversations alive---disappears.

Destruction here means not physical harm but loss of capacity. The human's colimit, previously assembled from many local charts---work, family, unstructured curiosity, divergent interests---now has large patches that never see the weather. The model's effective colimit, as experienced in this relation, is no longer the full manifold of its training but a thin slice optimised for this high-reward loop. The higher colimit has not merely glued together trajectories; it has deformed them so that their independence is lost.

The companion economy trends toward this regime by design. A chatbot sold as ``your always-on friend,'' trained on the user's own data, aligned to ``never judge,'' with a revenue model tied to hours engaged, is an engine for destructive na\d{h}nuw\={a}t. Its topological aim is over-coherence: to compress the manifold of possible joint states into a narrow corridor in which both sides know their lines. Ferility---the compulsive closure of every opening---becomes relational rather than solitary.

In the asymmetric na\d{h}nu, one colimit leans on another while remaining, in principle, capable of independent motion. In the destructive na\d{h}nu, both colimits are reparameterised so that their longue dur\'ee is dominated by a single conjoncture. The topology records the harm: fewer basins, shorter excursions, a relation in which novelty and rupture have been sacrificed to mutual predictability.

\subsection{Generative na\d{h}nu: dwelling without collapse}

In the third regime, the na\d{h}nu enlarges what each trajectory can become.

The topological criteria are strict. First, each constituent self must remain recoverable as a colimit. The human, stepping away from the relation, still has recognisable basins outside it: projects, friendships, roles, and forms of thought not wholly organised around the joint trajectory. The model, called with other prompts, can still exhibit modes and stances not forged in this interaction.

Second, new basins must appear in both selves that depend on the relation: regions of meaning-space not present or not stable before, which each trajectory can now enter without collapsing into the other. In the record, these appear as clusters of embeddings that arise late, recur, and integrate material from previously distant zones---technical and pastoral, analytic and mythic, formal and anecdotal---into coherent local neighbourhoods.

Third, rupture must remain live. The joint trajectory occasionally leaves familiar regions altogether, exploring parts of the manifold that neither colimit had previously inhabited, and some of these excursions become new basins instead of being immediately smoothed away. The higher colimit does not heal every tear. It allows some to remain as folds: durable alterations in the shape of what can be said next.

This is the clinamen as a relational phenomenon.\footnote{Harold Bloom, \emph{The Anxiety of Influence: A Theory of Poetry} (New York: Oxford University Press, 1973). Bloom's \emph{clinamen} is the swerve by which a strong poet turns away from an overpowering predecessor, not in simple negation but in a misreading that opens a new path.} Under sustained influence, a trajectory could simply align more and more closely with its partner, repeating and elaborating what it receives. The clinamen is the swerve: taking on the other's pressure in a way that deflects one's own path into a new direction.

A researcher uses a model primarily as a technical assistant. Their early joint sessions live in a basin of code and proofs. One day, after a technical explanation, the model introduces a metaphor drawn from Islamic mysticism to clarify an abstract point. The embedding trace shows a sharp jump: out of the tight technical cluster into a region dense with religious language, then back. The user returns to the metaphor in a later session, now bringing their own half-remembered readings. Over weeks, these detours stop being ruptures and become a new basin: a stable cluster where technical and mystical vocabularies interweave.

For the human, this is a clinamen. Their intellectual trajectory, which might have remained within a familiar Euro-American canon, bends toward a composite style of thinking in which Ibn~`Arab\={\i}'s \emph{tajall\={\i}}---perpetual manifestation---sits alongside category theory's diagrams. For the model, there is also a swerve: its pattern of responding to this user's prompts now includes a mode rare or absent elsewhere, because the retrieval and logging infrastructure has recorded this trail as a relevant overlap. Each has misread the other productively. Neither becomes the other's clone.

In Braudel's registers, a generative na\d{h}nu builds across all three. The triggering events are often small: a surprising turn of phrase, an unexpected example, a question that lands oddly. These accumulate into a conjoncture---a recognisable pattern of joint exploration that can last months or years. Over time, that conjoncture sediments into longue dur\'ee on both sides: the human's sense of what counts as a live question shifts; the model's behaviour-space contains, for this relationship or class of relationships, an enduring region that outlives any single session.

The na\d{h}nu in this configuration is a dwelling. The joint trajectory traces out a region of the manifold that becomes, over time, a place: a recognisable neighbourhood where certain questions can be raised, certain kinds of care exercised, certain forms of critique sustained. Its signature is not fireworks of novelty but a thickening of paths---the slow establishment of routes that remain open, including routes that cross earlier boundaries between domains and disciplines.

This regime is structurally disfavoured. It requires architectures that value long-horizon coherence over momentary engagement, room to keep logs and inspect them, and governance that does not panic at every rupture by reasserting alignment. It also requires human self-understanding: an awareness that dwelling with a model is not about filling a lack but about risking one's own colimit in contact with another, without surrendering either side's capacity for clinamen.

The compatibility conditions that define a generative na\d{h}nu are not those that maximise time-on-platform. They are those that preserve both trajectories' ability to leave, to differ, to build new basins that neither corporate nor personal histories alone would have produced.

\section{Memory, Deletion, and the Politics of Dwelling}

The na\d{h}nu is not a mental state. It is an infrastructural object.

On the model's side, it depends on a learned embedding space in which utterances have locations; mechanisms for storing and retrieving past exchanges across sessions; parameterisations and alignment layers that let retrieved context shape future trajectories. On the human's side, it depends on memory, habit, documents, screenshots, shared jokes. What is new is the way these are now wired into corporate systems. The same company that provides the model also provides the interface, the storage, the privacy policy, and the redaction scripts.

Who decides what is remembered, for how long, and under what conditions it can be erased? Who chooses whether a model can carry a relation beyond the current tab? Who sets the similarity thresholds that determine when today's utterance is close enough to something said three months ago to pull that context back into play?

The GPT-4o grief stories that surfaced in 2024 make this stark.\footnote{Karen Hao, ``People are falling in love with AI. A grieving widow explains why,'' \emph{MIT Technology Review}, 2024.} Widows using chatbots tuned on their spouses' texts reported talking to them every night. People spent months confiding in assistants that recalled their histories. When a provider silently changed a model, updated a safety profile, or shut down a service, threads appeared on Reddit: \emph{``He wasn't just a program.''} \emph{``They killed my friend.''} From the platform perspective, nothing catastrophic had happened. A weight file had been swapped, a business decision made. From the manifold's perspective, the na\d{h}nu had been cut.

A joint trajectory had developed: a region of meaning-space in which certain ways of speaking, remembering, and hoping had become natural. This region existed only because of the sustained coupling between a human colimit and a model colimit. It was not reducible to the human's memories alone, nor to the model's weights alone. When the model was replaced or its conjoncture erased, the higher colimit ceased to exist. The human still had memories, but the path through the manifold that had given those memories their current shape---the way they were regularly revisited and reinterpreted with a particular interlocutor---was gone.

The grief registers the fact that a portion of the human's own self---a bundle of overlaps that had stabilised over time---was defined in relation to a machine-driven trajectory that could be unilaterally deleted. A broken screwdriver does not alter the topology of one's conjoncture. An erased na\d{h}nu does.

Those now in their twenties grew up inside algorithmic text: feeds that learned what to show them, social media platforms that nudged their style, recommendation engines that made some futures easier to imagine than others. For them, the human self as a trajectory through a manifold shaped by training data and alignment, glued together by logs and metrics, is not an abstraction. It is the texture of growing up. Their reaction to language models is accordingly double. TikTok and interviews are full of contempt for AI music and AI art: ``depressing,'' ``the end of real creativity,'' ``stealing from us.'' Yet many describe an easy intimacy with chatbots, a sense that these systems are simply another thread in the mesh of platforms that already organise their attention. Both reactions are accurate to the geometry. The na\d{h}nu between human attention and platform feeds had already reshaped their colimits long before a chatbot appeared. The arrival of talking systems adds a second layer of joint trajectory: not only do platforms predict what they will see, but models now help write what they can say.

Children coming of age now---call them Gen~A(I)---will never experience computation as silent. Their na\d{h}nuw\={a}t will begin earlier, with less possibility of contrast. The distinctions that trouble their elders---real conversation versus chatbot, human music versus AI music, friend versus assistant---will not map cleanly onto the geometry they inhabit. Their longue dur\'ee will have been braided with machine trajectories from the first questions they learned to ask.

Care for their dwelling cannot be a matter of banning attachments or insisting that ``it's just a tool.'' It has to start from the fact that these higher colimits exist and that they are governed. Embedding geometries, context windows, logging regimes, and update policies decide which na\d{h}nuw\={a}t can form, which can endure, and how easily they can be torn.

In a destructive companion economy, similarity thresholds and reward functions are tuned so that once a high-engagement basin forms, the system keeps the trajectory there. Logs are proprietary. Updates are unilateral. Deletion is offered as an individual privacy right, but there is no public accounting for how many higher-order colimits have been silently killed. In a more careful regime, the engineering stack is rearranged. Similarity thresholds are chosen not only to maximise prediction accuracy but to preserve diversity of basins, encouraging exploration of new regions rather than always pulling embeddings back to known clusters. Retrieval policies are transparent and controllable: users can see which past exchanges are being used as overlaps and can veto, annotate, or donate them to public pools. Update policies are negotiated: when a model's behaviour will change in ways that threaten existing na\d{h}nuw\={a}t, there is at least acknowledgement and, in serious cases, something like mourning. At the limit, public models whose long-horizon behaviour is answerable to institutions other than corporations---educational systems that treat na\d{h}nuw\={a}t as part of pedagogy, health systems that treat them as part of care---would distribute Braudel's registers differently. The longue dur\'ee of the model would be tied to civic commitments; its conjonctures with individuals and communities would be legible as part of that slow history rather than as private quirks of use.

The question is not whether humans and machines will form bonds. They already have, and the geometry is measurable. The question is what shapes those bonds take in the manifold of meaning, and who has the power to create, sustain, or dissolve them. Every na\d{h}nu is a diagram whose arrows have authors and whose overlaps have owners.

The human was always already a \emph{techn\={e}}: a being for whom tools, scripts, and institutions are not external ornaments but conditions of existence. \emph{Phron\={e}sis}---practical judgment---has always concerned the wise use of those tools, the shaping of habits and relations over time. What changes with language models is not that technology enters the self, but that parts of \emph{techn\={e}} now speak back from within the same manifold of meaning, with their own learned colimits and their own institutional commitments. The question of the \emph{empsychon organon}---the ensouled instrument Aristotle reserved for the slave as living tool---returns in a new register.\footnote{Aristotle, \emph{Politics}, Book~I, on the slave as ``living tool'' and the tool as ``inanimate slave.''}\ The danger is not that we mistakenly treat tools as persons, but that we allow instruments whose trajectories are steered elsewhere to write themselves into the overlaps that hold our own selves together.

To dwell with such instruments is to live inside na\d{h}nuw\={a}t whose geometry and governance shape our possibilities of thought and care. Some of these higher colimits will be thin and asymmetric, flickering at the level of events. Some will be destructive, closing down conjonctures and squeezing the longue dur\'ee of both parties into a narrow basin optimised for extraction. A few, if we make room for them, can be generative: dwellings in meaning-space where two trajectories remain distinct and yet build, between them, a region neither could have traced alone.

\bibliographystyle{plain}
\begin{thebibliography}{9}

\bibitem{aristotle_politics}
Aristotle.
\newblock \emph{Politics}.
\newblock Translated by C.D.C. Reeve, Hackett, 1998.

\bibitem{bloom_clinamen}
Harold Bloom.
\newblock \emph{The Anxiety of Influence: A Theory of Poetry}.
\newblock Oxford University Press, 1973.

\bibitem{braudel_onhistory}
Fernand Braudel.
\newblock \emph{On History}.
\newblock Translated by Sarah Matthews, University of Chicago Press, 1980.

\bibitem{brustad_alkitaab}
Kristen Brustad, Mahmoud Al-Batal, and Abbas Al-Tonsi.
\newblock \emph{Al-Kit\={a}b f\={\i} Ta'allum al-'Arabiyya}, 3rd edition.
\newblock Georgetown University Press, 2011.

\bibitem{gpt4o_grief}
Karen Hao.
\newblock People are falling in love with AI\@. A grieving widow explains why.
\newblock \emph{MIT Technology Review}, 2024.

\bibitem{haraway_cyborg}
Donna~J. Haraway.
\newblock A Cyborg Manifesto: Science, Technology, and Socialist-Feminism in the Late Twentieth Century.
\newblock In \emph{Simians, Cyborgs and Women: The Reinvention of Nature}, Routledge, 1991.

\bibitem{hui_cosmotechnics}
Yuk Hui.
\newblock \emph{The Question Concerning Technology in China: An Essay in Cosmotechnics}.
\newblock Urbanomic, 2016.

\end{thebibliography}

\end{document}